% =====================================================================
% KTU Bakalauro baigiamojo projekto šablonas — 2026 m. versija
% Programų sistemos (6121BX012), Informatikos fakultetas
% =====================================================================

\documentclass[12pt]{article}

\usepackage{ktu_thesis_style}
\usepackage{lipsum} % Tik pavyzdiniam tekstui — pašalinkite rašydami savo darbą.

\begin{document}
% =====================================================================
% Darbo metaduomenys — KTU IF Programų sistemos (6121BX012)
% Bakalauro baigiamasis projektas — 2026 m. šablonas
% =====================================================================
%
% Pakeiskite žemiau esančias pavyzdines reikšmes savomis. Laukai
% suskirstyti į lietuviškus (naudojami viršeliuose ir darbo tekste) ir
% angliškus (naudojami tik santraukoje anglų kalba).
%
% Vardai rašomi „Pavardė Vardas“ tvarka, kaip KTU 2026 šablone.

% --------- Bendri ---------
\def\projectYear{2026}
\def\projectCity{Kaunas}

% --------- Lietuviškai (pagrindinis tekstas) ---------
\def\projectTitleLT{Darbo pavadinimas}
\def\projectTypeLT{Baigiamasis bakalauro studijų projektas}

\def\projectFacultyLT{Informatikos fakultetas}
\def\projectStudiesLT{Programų sistemos (6121BX012)}
\def\projectStudyFieldAndAreaLT{Informatikos mokslai, Programų sistemos}

\def\projectAuthorNameLT{Pavardė Vardas}
\def\projectSupervisorPositionLT{doc. dr.}
\def\projectSupervisorNameLT{Pavardė Vardas}
% Recenzentas paskiriamas studijų programos komiteto prieš gynimą.
% Jei recenzento dar nėra, palikite tuščią: \def\projectReviewerNameLT{}
% (tada antrajame viršelyje recenzento eilutė nerodoma).
\def\projectReviewerPositionLT{doc. dr.}
\def\projectReviewerNameLT{Pavardė Vardas}

\def\projectKeywordsLT{raktinis žodis, raktinis žodis, raktinis žodis}

% --------- Angliškai (tik santrauka) ---------
\def\projectTitleEN{Thesis Title}
\def\projectTypeEN{Bachelor's Final Degree Project}

\def\projectFacultyEN{Faculty of Informatics}
\def\projectStudyFieldAndAreaEN{Computer Sciences, Software Systems}

\def\projectAuthorNameEN{Name Surname}
\def\projectSupervisorPositionEN{assoc. prof. dr.}
\def\projectSupervisorNameEN{Name Surname}

\def\projectKeywordsEN{keyword, keyword, keyword}


% Puslapiai numeruojami nuo titulinio (counter=1), bet matomas numeris
% atsiranda tik nuo turinio (titulinio puslapio aplinka pati nustato
% \thispagestyle{empty}, taigi pirmieji puslapiai be numerio).
\pagenumbering{arabic}

% =====================================================================
% Pirmasis viršelis (announcement) — KTU 2026 sablonas
% Be parasu, vienas vadovas, be recenzento.
% =====================================================================

\definecolor{lines_name}{RGB}{212,175,54}

\begin{titlepage}
   \begin{center}

       \vspace*{0.3cm}
       \includegraphics[width=2.46cm]{images/ktu-ikona.pdf}
       \vspace{0.3cm}

       \textbf{Kauno technologijos universitetas}\\
       \projectFacultyLT

       \vspace{2.5cm}

       \LARGE
       \textbf{\projectTitleLT}\\[0.5em]
       \large
       \projectTypeLT

       \vspace{2cm}
       {\color{lines_name}\rule{9cm}{0.5pt}}
       \vspace{0.4cm}

       \normalsize
       \textbf{\projectAuthorNameLT}\\
       Projekto autorius

       \vspace{0.8cm}

       \textbf{\projectSupervisorPositionLT~\projectSupervisorNameLT}\\
       Vadovas

       \vspace{0.4cm}
       {\color{lines_name}\rule{9cm}{0.5pt}}

       \vfill

       \normalsize\textbf{\projectCity, \projectYear}
   \end{center}
\end{titlepage}

% =====================================================================
% Antrasis viršelis — formalus titulinis su parasu/data eilutemis,
% studiju programos kodu, vadovu ir recenzentu (KTU 2026 sablonas).
% =====================================================================

\definecolor{lines_name}{RGB}{212,175,54}

\begin{titlepage}
   \begin{center}

       \vspace*{0.3cm}
       \includegraphics[width=2.46cm]{images/ktu-ikona.pdf}
       \vspace{0.3cm}

       \textbf{Kauno technologijos universitetas}\\
       \projectFacultyLT

       \vspace{2cm}

       \LARGE
       \textbf{\projectTitleLT}\\[0.5em]
       \large
       \projectTypeLT\\
       \projectStudiesLT

       \vspace{2cm}

       \makebox[\linewidth][r]{%
       \begin{minipage}{9cm}
       \setlength{\parindent}{0pt}
       {\color{lines_name}\rule{\linewidth}{0.5pt}}\par
       \vspace{0.4cm}

       \normalsize
       \textbf{\projectAuthorNameLT}\\
       Projekto autorius

       \vspace{0.8cm}

       \textbf{\projectSupervisorPositionLT~\projectSupervisorNameLT}\\
       Vadovas

       \ifx\projectReviewerNameLT\empty\else
       \vspace{0.8cm}

       \textbf{\projectReviewerPositionLT~\projectReviewerNameLT}\\
       Recenzentas
       \fi

       \vspace{0.4cm}
       {\color{lines_name}\rule{\linewidth}{0.5pt}}
       \end{minipage}%
       }

       \vfill
       \normalsize\textbf{\projectCity, \projectYear}
   \end{center}
\end{titlepage}

% =====================================================================
% Akademinio sąžiningumo deklaracija — KTU 2026 šablonas (lietuviškai).
% =====================================================================

\definecolor{lines_name}{RGB}{212,175,54}

\begin{titlepage}
   \begin{center}

       \vspace*{0.3cm}
       \includegraphics[width=2.46cm]{images/ktu-ikona.pdf}
       \vspace{0.3cm}

       \textbf{Kauno technologijos universitetas}\\
       \projectFacultyLT\\
       \projectAuthorNameLT

       \vspace{2cm}

       \LARGE
       \textbf{\projectTitleLT}\\[0.5em]
       \large
       Akademinio sąžiningumo deklaracija

       \vspace{1cm}
   \end{center}

   \normalsize
   Patvirtinu, kad:
   \begin{enumerate}
     \item baigiamąjį projektą parengiau savarankiškai ir sąžiningai, nepažeisdama(s) kitų asmenų autoriaus ar kitų teisių, laikydamasi(s) Lietuvos Respublikos autorių teisių ir gretutinių teisių įstatymo nuostatų, Kauno technologijos universiteto (toliau – Universitetas) intelektinės nuosavybės valdymo ir perdavimo nuostatų bei Universiteto akademinės etikos kodekse nustatytų etikos reikalavimų;
     \item baigiamajame projekte visi pateikti duomenys ir tyrimų rezultatai yra teisingi ir gauti teisėtai, nei viena šio projekto dalis nėra plagijuota nuo jokių spausdintinių ar elektroninių šaltinių, visos baigiamojo projekto tekste pateiktos citatos ir nuorodos yra nurodytos literatūros sąraše;
     \item įstatymų nenumatytų piniginių sumų už baigiamąjį projektą ar jo dalis niekam nesu mokėjęs (-usi);
     \item suprantu, kad išaiškėjus nesąžiningumo ar kitų asmenų teisių pažeidimo faktui, man bus taikomos akademinės nuobaudos pagal Universitete galiojančią tvarką ir būsiu pašalinta(s) iš Universiteto, o baigiamasis projektas gali būti pateiktas Akademinės etikos ir procedūrų kontrolieriaus tarnybai nagrinėjant galimą akademinės etikos pažeidimą.
   \end{enumerate}

   \begin{flushright}
       \projectAuthorNameLT\\
       \textit{Patvirtinta elektroniniu būdu}
   \end{flushright}

   \vfill
\end{titlepage}


% --- Santrauka LT ir EN ---
% =====================================================================
% Santrauka — KTU 2026 šablonas.
% =====================================================================

\begin{titlepage}

\projectAuthorNameLT. \projectTitleLT. \projectTypeLT\ / vadovas \projectSupervisorPositionLT~\projectSupervisorNameLT; Kauno technologijos universitetas, \projectFacultyLT.

Studijų kryptis ir sritis (studijų krypčių grupė): \projectStudyFieldAndAreaLT.

Reikšminiai žodžiai: \projectKeywordsLT.

\projectCity, \projectYear. \pageref{LastPage}~p.

\vspace{0.5cm}
\begin{center}
  \textbf{Santrauka}
\end{center}

Santraukoje glaustai aprašoma darbo problema, objektas, tikslas, uždaviniai, taikyti metodai, pagrindiniai rezultatai ir išvados. Rekomenduojama apimtis -- iki vieno puslapio.

\lipsum[1-2]

\end{titlepage}

% =====================================================================
% Summary — KTU 2026 template (English).
% =====================================================================

\begin{otherlanguage}{english}
\begin{titlepage}

\projectAuthorNameEN. \projectTitleEN. \projectTypeEN\ / supervisor \projectSupervisorPositionEN\ \projectSupervisorNameEN; \projectFacultyEN, Kaunas University of Technology.

Study field and area (study field group): \projectStudyFieldAndAreaEN.

Keywords: \projectKeywordsEN.

\projectCity, \projectYear. \pageref{LastPage}~pages.

\vspace{0.5cm}
\begin{center}
  \textbf{Summary}
\end{center}

The summary briefly describes the problem, object, aim, objectives, methods, main results and conclusions of the project. Recommended length -- up to one page.

\lipsum[1-2]

\end{titlepage}
\end{otherlanguage}


% --- Turinys ---
\clearpage
\tableofcontents

% --- Lentelių sąrašas ---
\clearpage
\phantomsection
\addcontentsline{toc}{section}{\listtablename}
\listoftables

% --- Paveikslų sąrašas ---
\clearpage
\phantomsection
\addcontentsline{toc}{section}{\listfigurename}
\listoffigures

% --- Santrumpų ir terminų sąrašas ---
\clearpage
\phantomsection
\addcontentsline{toc}{section}{Santrumpų ir terminų sąrašas}
\section*{Santrumpų ir terminų sąrašas}
\textbf{Santrumpos}

\begin{description}[leftmargin=0pt,labelsep=0.5em,itemsep=2pt,parsep=0pt]
\item[API] -- programų sąsaja (angl. \textit{Application Programming Interface}); sutartinis būdas, kuriuo programos arba jų moduliai bendrauja tarpusavyje.
\item[DI] -- dirbtinis intelektas (angl. \textit{Artificial Intelligence}).
\item[PA] -- panaudojimo atvejis (angl. \textit{Use Case}); UML naudotojo elgsenos diagramos elementas.
\item[UML] -- unifikuota modeliavimo kalba (angl. \textit{Unified Modeling Language}).
\end{description}

\textbf{Terminai}

\begin{description}[leftmargin=0pt,labelsep=0.5em,itemsep=2pt,parsep=0pt]
\item[Terminas] -- termino apibrėžimas.
\item[Kitas terminas] -- termino apibrėžimas.
\end{description}


% --- Įvadas ---
\clearpage
\phantomsection
\addcontentsline{toc}{section}{Įvadas}
\section*{Įvadas}
Įvade pagrindžiamas temos aktualumas, suformuluojama problema, darbo objektas, tikslas ir uždaviniai. Uždaviniai turi būti tiesiogiai atspindėti skyriuje „Rezultatai ir išvados“.

\lipsum[1]

Darbo tikslas -- \lipsum[2][1-2]

Šiam tikslui pasiekti išsikelti šie uždaviniai:

\begin{enumerate}
    \item Išanalizuoti \ldots;
    \item Suformuluoti reikalavimus \ldots;
    \item Suprojektuoti ir įgyvendinti \ldots;
    \item Įvertinti sukurtos sistemos kokybę \ldots.
\end{enumerate}

Darbą sudaro penki numeruoti skyriai ir nenumeruotas baigiamasis skyrius. \lipsum[3][1-3]


% --- Pagrindiniai skyriai (numeruoti) ---
\section{Sistemos analizė}

Šiame skyriuje pateikiamas techninis pasiūlymas ir sistemos techninių galimybių analizė.

\subsection{Techninis pasiūlymas}

\subsubsection{Sistemos apibrėžimas}

\lipsum[1]

\subsubsection{Pagrindinis sistemos tikslas}

\lipsum[2]

\subsubsection{Sistemos pagrįstumas}

Teiginiai iš literatūros visada pagrindžiami nuoroda į šaltinį \citep{vysniauskas2006}. Kelios nuorodos sujungiamos automatiškai \citep{masiulis2007, birziska1929, gradauskas2000}, elektroniniai šaltiniai cituojami taip pat \citep{litnet2012, valiulyte2000, ieee2009}.

\lipsum[3]

\subsubsection{Egzistuojančių sprendimų analizė}

Lentelės pavadinimas rašomas virš lentelės, o tekste į lentelę būtina paminėti nuorodą prieš pačią lentelę (žr. \ref{tab:palyginimas} lentelę). Numeris ir žodis „lentelė.“ pridedami automatiškai.

\begin{table}[!ht]
\caption{Egzistuojančių sprendimų palyginimas}
\label{tab:palyginimas}
\centering
\begin{tabular}{|p{4cm}|p{3.5cm}|p{3.5cm}|p{3.5cm}|} \toprule
\textbf{Kriterijus} & \textbf{Sprendimas A} & \textbf{Sprendimas B} & \textbf{Kuriama sistema} \\ \midrule
Kriterijus 1 & Taip & Ne & Taip \\ \hline
Kriterijus 2 & Ne & Iš dalies & Taip \\ \hline
Kriterijus 3 & Taip & Taip & Taip \\ \bottomrule
\end{tabular}
\end{table}

\subsubsection{Prototipai ir pagalbinė informacija}

Paveikslo pavadinimas rašomas po paveikslu. Nuoroda į paveikslą pateikiama tekste prieš jį (žr. \ref{fig:pavyzdys} pav.).

\begin{figure}[!ht]
\centering
\includegraphics[width=3cm]{images/ktu-ikona.pdf}
\caption{Pavyzdinis paveikslas}
\label{fig:pavyzdys}
\end{figure}

\subsection{Sistemos techninių galimybių analizė}

\subsubsection{Technologijų apžvalga}

Formulės numeruojamos pagal skyrių:
\begin{equation}
  \label{eq:pavyzdys}
  P = \frac{n_{\text{padengta}}}{n_{\text{iš viso}}} \cdot 100~\%
\end{equation}
Formulėje (\ref{eq:pavyzdys}) $P$ -- padengimo procentas.

\subsubsection{Pasirinktas techninis sprendimas}

\lipsum[4]

\section{Sistemos kūrimo projektas}

Šiame skyriuje pateikiama sistemos reikalavimų specifikacija, projektavimo metodai, statinis ir dinaminis sistemos vaizdas.

\subsection{Sistemos reikalavimų specifikacija}

\subsubsection{Komercinė specifikacija}

\lipsum[1]

\subsubsection{Sistemos funkcijos}

Funkciniai reikalavimai išreiškiami panaudojimo atvejų (toliau -- PA) modeliu. Kiekvienas PA aprašomas specifikacijos lentele, kurią sukuria makrokomanda \verb|\usecasespec| (žr. \ref{tab:pa-perziureti-irasus} lentelę). Reikšmes, kuriose yra kablelių, būtina apgaubti riestiniais skliaustais.

\usecasespec{
  label=tab:pa-perziureti-irasus,
  pa=Peržiūrėti įrašus,
  tikslas=Atvaizduoti visus naudotojui prieinamus įrašus.,
  aprasymas={Prisijungęs naudotojas pasirenka meniu punktą „Įrašai“. Sistema pateikia įrašų sąrašą, kurį galima filtruoti ir rikiuoti.},
  pries=Naudotojas prisijungęs prie sistemos.,
  aktoriai={Naudotojas, Administratorius},
  ispleciantys=Peržiūrėti įrašo informaciją,
  apimami=--,
  specializuojami=--,
  po=Naudotojas mato įrašų sąrašą.
}

\subsubsection{Naudotojo sąsajos specifikacija}

\lipsum[2]

\subsubsection{Realizacijai keliami reikalavimai}

\lipsum[3]

\subsubsection{Techninė specifikacija}

\lipsum[4]

\subsection{Sistemos projektavimo metodai}

\subsubsection{Projektavimo valdymas ir eiga}

\lipsum[5]

\subsubsection{Projektavimo technologija}

\lipsum[6]

\subsubsection{Programavimo kalbos, derinimo, automatizavimo priemonės, operacinės sistemos}

\lipsum[7]

\subsection{Sistemos projektas}

\subsubsection{Statinis sistemos vaizdas}

\paragraph{Diegimo diagrama.} \lipsum[8][1-4]

\paragraph{Komponentų diagrama.} \lipsum[9][1-4]

\subsubsection{Dinaminis sistemos vaizdas}

\lipsum[10]

\section{Sistemos testavimas}

\subsection{Testavimo planas ir kriterijai}

\lipsum[1]

\subsection{Vienetų ir integracinis testavimas}

\lipsum[2]

\subsection{Kodo padengimas}

Skaičiai rašomi su dešimtainiu kableliu ir nedalomu tarpu prieš vienetą, pvz., 83{,}15~\%, 500~ms, 17\,047 eilutės.

\subsection{Statinė kodo analizė}

\lipsum[3]

\section{Sistemos dokumentacija naudotojui}

\subsection{Apibendrintas sistemos galimybių aprašymas}

\lipsum[1]

\subsection{Naudotojo vadovas}

\lipsum[2]

\subsection{Diegimo vadovas}

\lipsum[3]

\subsection{Administratoriaus vadovas}

\lipsum[4]

\section{Dirbtinio intelekto naudojimo savirefleksija}

Pagal KTU 2026 m. šabloną šis skyrius yra privalomas: jame aprašoma, kokie dirbtinio intelekto (DI) įrankiai naudoti, kuriuose darbo etapuose ir kaip buvo tikrinami jų rezultatai.

\subsection{Naudoti DI įrankiai ir kontekstas}

\lipsum[1]

\subsection{Taikymas projekto kūrimo etapuose}

Apibendrinta DI naudojimo apžvalga pateikiama \ref{tab:di-etapai} lentelėje.

\begin{table}[!ht]
\caption{DI panaudojimas projekto kūrimo etapuose}
\label{tab:di-etapai}
\centering
\begin{tabular}{|p{2.6cm}|p{6.8cm}|p{5.8cm}|} \toprule
\textbf{Etapas} & \textbf{DI veikla} & \textbf{Verifikavimo metodas} \\ \midrule
Analizė & \ldots & \ldots \\ \hline
Projektavimas & \ldots & \ldots \\ \hline
Programavimas & \ldots & \ldots \\ \hline
Testavimas & \ldots & \ldots \\ \hline
Dokumentavimas & \ldots & \ldots \\ \bottomrule
\end{tabular}
\end{table}

\subsection{Stiprybės ir trūkumai}

\lipsum[2]


% --- Rezultatai ir išvados ---
\clearpage
\phantomsection
\addcontentsline{toc}{section}{Rezultatai ir išvados}
\section*{Rezultatai ir išvados}
Išvados pateikiamos ta pačia tvarka kaip įvade suformuluoti uždaviniai -- kiekvienam uždaviniui po vieną išvadą.

\begin{enumerate}
    \item \lipsum[1][1-3]
    \item \lipsum[2][1-3]
    \item \lipsum[3][1-3]
    \item \lipsum[4][1-3]
\end{enumerate}


% --- Literatūros sąrašas ---
\clearpage
\phantomsection
\addcontentsline{toc}{section}{Literatūros sąrašas}
\renewcommand{\refname}{Literatūros sąrašas}
\bibliography{references}

% --- Priedai ---
\clearpage
\phantomsection
\addcontentsline{toc}{section}{Priedai}
\section*{Priedai}
\setcounter{table}{0}
\setcounter{figure}{0}
\renewcommand{\thetable}{P\arabic{table}}
\renewcommand{\thefigure}{P\arabic{figure}}
% Priedų lentelės/paveikslai nepatenka į Lentelių/Paveikslų sąrašus.
\captionsetup{list=no}
Kiekvienas priedas pradedamas makrokomanda \verb|\priedas{Pavadinimas}|, kuri sukuria antraštę „N priedas. Pavadinimas“ ir įtraukia ją į turinį. Priedų lentelės ir paveikslai numeruojami „P1“, „P2“ ir t.\,t.

\priedas{Pirmojo priedo pavadinimas}

\lipsum[1]

\begin{table}[!ht]
\caption{Pavyzdinė priedo lentelė}
\label{tab:priedas-pavyzdys}
\centering
\begin{tabular}{|p{6cm}|r|} \toprule
\textbf{Rodiklis} & \textbf{Reikšmė} \\ \midrule
Rodiklis 1 & 100{,}0 \\ \hline
Rodiklis 2 & 42{,}0 \\ \bottomrule
\end{tabular}
\end{table}

\priedas{Antrojo priedo pavadinimas}

\lipsum[2]


\clearpage
\label{LastPage}
\end{document}
